% Generated from article.tex and its local inputs. Edit modular sources.
\documentclass[11pt,letterpaper]{article}
\usepackage[margin=1in,headheight=15pt]{geometry}
\usepackage{amsmath,amssymb,amsthm,mathtools,bm}
\usepackage{lmodern,microtype}
\usepackage[T1]{fontenc}
\usepackage{booktabs,array,longtable,enumitem}
\usepackage{xcolor}
\definecolor{ink}{HTML}{15344A}
\usepackage[colorlinks=true,linkcolor=ink,citecolor=ink,urlcolor=ink,breaklinks=true]{hyperref}
\usepackage{fancyhdr}
\pagestyle{fancy}\fancyhf{}
\fancyhead[L]{\small Determinantal harmonic resonance}
\fancyhead[R]{\small ProveIt research continuation}
\fancyfoot[C]{\thepage}
\renewcommand{\headrulewidth}{0.3pt}
\setlength{\parskip}{0.35em}
\setlength{\parindent}{1.25em}
\setlist{itemsep=0.3em,topsep=0.4em}
\numberwithin{equation}{section}
\newtheorem{theorem}{Theorem}[section]
\newtheorem{lemma}[theorem]{Lemma}
\newtheorem{proposition}[theorem]{Proposition}
\newtheorem{corollary}[theorem]{Corollary}
\theoremstyle{definition}\newtheorem{definition}[theorem]{Definition}
\theoremstyle{remark}\newtheorem{remark}[theorem]{Remark}
\DeclareMathOperator{\Li}{Li}
\DeclareMathOperator{\FP}{FP}
\DeclareMathOperator{\Res}{Res}
\DeclareMathOperator{\Rea}{Re}
\newcommand{\C}{\mathbb C}\newcommand{\Q}{\mathbb Q}\newcommand{\N}{\mathbb N}
\newcommand{\eps}{\varepsilon}
\newcommand{\fall}[2]{#1^{\underline{#2}}}
\newcommand{\diff}{\,\mathrm d}
\newcommand{\ind}{\mathbf1}
\newcommand{\repo}{\texttt{ProveIt}}
\newcommand{\Bell}{\mathcal B}
\hypersetup{pdftitle={Determinantal Harmonic Resonance},pdfauthor={OpenAI-assisted research draft for the ProveIt project},pdfsubject={Gauss--Vandermonde finite parts, Stieltjes jets, and depth-two polylogarithmic closure}}
\begin{document}
\begin{titlepage}
\centering
\vspace*{1.1cm}
{\small\scshape Analysis / Polylogarithms\par}
\vspace{1cm}
{\huge\bfseries\color{ink} Determinantal Harmonic\\[0.2em] Resonance\par}
\vspace{0.7cm}
{\Large Gauss--Vandermonde finite parts,\\[0.25em] Stieltjes jets, and depth-two\\[0.25em] polylogarithmic closure\par}
\vspace{1cm}
{\large A research continuation for the ProveIt project\par}
\vspace{0.25cm}
{OpenAI-assisted mathematical research draft\par}
{10 October 2026\par}
\vspace{1cm}
\begin{minipage}{0.91\textwidth}
\small
\textbf{Abstract.}
We develop a determinantal extension of the Gamma--harmonic identity
programme. A self-contained evaluation of a discrete Gauss--Vandermonde
sum is used at its first logarithmic resonance. Separating the largest
summation index produces a normally convergent subtraction, an explicit
hard-cutoff constant, and all-order identities coupling many-variable
harmonic sums to generalized Stieltjes constants. The same construction
has exact Hurwitz-zeta derivatives and anchored log-Gamma primitives.

For rational weights with simple poles on the negative integer grid,
we prove that the finite moment matrix has harmonic dependence in only
one entry. Consequently an arbitrary number of summation variables still
admits a depth-two polylogarithmic generating function. Every positive
Euler primitive at the endpoint reduces to ordinary zeta values and
products of two zeta values. We give explicit two-, three-, and four-variable
evaluations, an all-order mixed harmonic generating function, and a
common-translate cyclotomic extension.

The proofs, exact rational reducers, and diagnostic calculations are
separate parts of the delivery. The replay contains 732 exact assertions
and 85 numerical diagnostics at 100-digit working precision. The latter
are not interval certificates. We do not assert global priority, resolve
the manuscript's $S_6$ or $S_8$ conjectures, or confer canonical or
proof-assistant status on this draft.
\end{minipage}
\vfill
{\small Source anchor: \texttt{87c7db629552afbc2e803001526debe3430a31b0}\par}
{\small Scope and integration qualifications are recorded in Section~\ref{dhr:sec:audit}.\par}
\end{titlepage}
\begingroup
\small
\setlength{\parskip}{0pt}
\tableofcontents
\endgroup
\clearpage
% BEGIN sections/01-context.tex
\section{Context, scope, and conventions}
\label{dhr:sec:context}

\subsection{The continuation problem}
The canonical \repo\ manuscript already has a substantial Gamma--harmonic
and Hurwitz--Stieltjes calculus. In particular, the inspected mixed-Dougall
chapter proves a normally convergent two-parameter Gamma germ, all its
fixed mixed derivatives, a weighted trigamma square, and a quadratic
zeta-tail family \cite{repo:readme,repo:dougall}. The manuscript also records
an all-order Gauss proof of the corrected pure and mixed Choi harmonic
families. Re-proving the one-variable coefficient engine would not by
itself be a meaningful continuation.

We instead ask what happens when a squared Vandermonde couples an arbitrary
number of summation indices. There are two distinct questions. First,
can the Gamma product still supply a correctly normalized Stieltjes
subtraction when the sum ceases to converge? Second, does the increase
in the number of indices force a comparable increase in polylogarithmic
depth? We answer both for the first resonant determinantal family.

The first answer requires attention to a distinction easily hidden by
formal Laurent expansion: a spectral finite part is generally \emph{not}
the constant term of a hard cutoff. The residue amplitude varies with
the spectral parameter. Its derivative is precisely the discrepancy.
The second answer is more economical than the number of variables
suggests: for simple-pole rational weights at critical decay, only the
bottom-right moment contains a harmonic number. This gives depth at most
two at every particle number.

\subsection{Main results and what they do not claim}
Theorem~\ref{dhr:thm:gauss} evaluates the convergent discrete partition
function by elementary determinant algebra and Gauss summation. Its
classical ingredients are not presented as new. Theorems
\ref{dhr:thm:resonance} and \ref{dhr:thm:jets} give the normally convergent
resonance calculus and its all-order Stieltjes consequences. Theorems
\ref{dhr:thm:rank-one} and \ref{dhr:thm:closure} give the finite-cutoff
rank-one harmonic theorem and its all-particle depth-two consequence.
The concrete formulas in Section~\ref{dhr:sec:explicit} are exact
applications, not identities inferred from a numerical relation search.

These are results proved in the present report. An exhaustive literature
comparison has not established their global novelty. Nor do they prove
a lower bound on the depth of any actual period. The manuscript explicitly
distinguishes formal cyclotomic quotients from numerical period
independence \cite{repo:quotients}; we retain that distinction. The short
$S_6$ and $S_8$ identities remain outside the results proved here.

\subsection{Notation and analytic conventions}
For $k\geq0$, write $(x)_k=\Gamma(x+k)/\Gamma(x)$ for a rising
factorial, and $\fall{x}{k}=x(x-1)\cdots(x-k+1)$ for a falling factorial.
Empty products and empty determinants are one. We use
\[
 H_n^{(s)}=\sum_{j=1}^n j^{-s},\qquad H_n=H_n^{(1)},\qquad H_0^{(s)}=0.
\]
The gamma logarithmic derivative is $\psi=\Gamma'/\Gamma$. Our
Stieltjes convention is
\begin{equation}\label{dhr:eq:stieltjes-convention}
 \zeta(1+u,\alpha)=\frac1u+
 \sum_{k\geq0}\frac{(-1)^k}{k!}\gamma_k(\alpha)u^k,
 \qquad \Rea\alpha>0.
\end{equation}
Thus $\gamma_0(\alpha)=-\psi(\alpha)$ and $\gamma_0(1)=\gamma$.
Derivatives $\zeta^{(k)}(s,\alpha)$ are in the first, or order,
variable unless specified otherwise. We use the standard branches of
Gamma and Hurwitz zeta on the stated right half-planes
\cite{dlmf:gamma,dlmf:hurwitz,dlmf:polygamma}.

Within $|z|<1$, define
\begin{equation}\label{dhr:eq:polylogs}
 \Li_s(z)=\sum_{n\geq1}\frac{z^n}{n^s},\qquad
 \Li_{p,1}(z,1)=\sum_{n>m\geq1}\frac{z^n}{n^pm}.
\end{equation}
The latter convention agrees with the two-argument convention
$\Li_{p,1}(x,y)=\sum_{n>m\geq1}x^ny^m/(n^pm)$.
In particular $\Li_{1,1}(z,1)=\tfrac12\log^2(1-z)$. All formulas
initially in the disk use the logarithm analytic there and zero at
$z=0$. Rational Laurent factors in $z$ in the displayed reductions
have removable singularities after the full expression is assembled.
Analytic continuation means continuation of this common germ, not an
independent choice of branch for each summand.

The complete exponential Bell polynomials are normalized by
\begin{equation}\label{dhr:eq:bell}
 \exp\left(\sum_{j\geq1}x_j\frac{t^j}{j!}\right)
 =\sum_{k\geq0}\Bell_k(x_1,\ldots,x_k)\frac{t^k}{k!}.
\end{equation}
For example, $\Bell_1=x_1$, $\Bell_2=x_1^2+x_2$, and
$\Bell_3=x_1^3+3x_1x_2+x_3$.

Whenever a series of subtracted brackets is asserted to converge,
the brackets are the summands. This does not license splitting the
series into separately divergent components. All index and parameter
derivatives in the results below are fixed finite orders.

% END sections/01-context.tex

% BEGIN sections/02-determinant.tex
\section{A discrete Gauss--Vandermonde product}
\label{dhr:sec:determinant}

For $r\geq1$ and $0\leq n_1<\cdots<n_r$, put
\[
 \Delta(\bm n)=\prod_{1\leq i<j\leq r}(n_j-n_i),\qquad
 w_n(a,c;b)=\frac{\Gamma(n+a)\Gamma(n+c)}
 {\Gamma(n+b)\Gamma(n+1)}.
\]
Define
\begin{equation}\label{dhr:eq:partition}
 Z_r(a,c;b)=\sum_{0\leq n_1<\cdots<n_r}
       \Delta(\bm n)^2\prod_{i=1}^r w_{n_i}(a,c;b),
 \qquad Z_0=1.
\end{equation}
The ordering is part of the definition. There is no additional $r!$
in any product formula below.

\begin{theorem}[Gauss--Vandermonde product]\label{dhr:thm:gauss}
Let $\Rea a,\Rea c>0$ and $\Rea d>2r-2$, where $d=b-a-c$.
Then \eqref{dhr:eq:partition} converges absolutely and locally normally,
and
\begin{equation}\label{dhr:eq:gauss-product}
 \boxed{\displaystyle
 Z_r(a,c;b)=\prod_{j=0}^{r-1}
 \frac{j!\,\Gamma(a+j)\Gamma(c+j)\Gamma(d-r+1-j)}
 {\Gamma(b-a-j)\Gamma(b-c-j)}.}
\end{equation}
For positive real parameters in this domain all summands are positive.
\end{theorem}

We give the determinant argument in detail, both to make the normalizing
factor checkable and to avoid relying on an unproved discrete-Selberg
analogy. Gamma ratio estimates and Gauss summation are classical
\cite{dlmf:gamma,dlmf:gauss}; the determinant technique belongs to the
standard evaluation toolkit \cite{krattenthaler}.

\begin{lemma}[A rational factorial determinant]\label{dhr:lem:rational-det}
For indeterminates $x_0,\ldots,x_{r-1},t$,
\begin{equation}\label{dhr:eq:rational-det}
 \det_{0\leq i,j<r}\frac{(x_i)_j}{(x_i+t)_j}
 =\frac{\displaystyle\prod_{j=0}^{r-1}(t)_j
                   \prod_{0\leq i<k<r}(x_k-x_i)}
        {\displaystyle\prod_{i=0}^{r-1}(x_i+t)_{r-1}}.
\end{equation}
The identity holds at every specialization where the displayed
rational functions are defined, and elsewhere by cancellation.
\end{lemma}
\begin{proof}
Multiply row $i$ by $(x_i+t)_{r-1}$. The resulting column polynomials
are
\[
 p_j(x)=(x)_j(x+t+j)_{r-1-j},\qquad 0\leq j<r.
\]
They have degree at most $r-1$, so their evaluation determinant is
the Vandermonde in the $x_i$ times their coefficient determinant.
To compute that coefficient determinant, evaluate at $x_i=-i$.
The evaluation matrix is lower triangular, since $(-i)_j=0$ for
$j>i$, with diagonal
\[
 p_i(-i)=(-1)^i i!\,(t)_{r-1-i}.
\]
Its determinant divided by
$\prod_{i<k}((-k)-(-i))=(-1)^{r(r-1)/2}\prod_{k=0}^{r-1}k!$
is $\prod_{j=0}^{r-1}(t)_j$. Dividing the row factors proves
\eqref{dhr:eq:rational-det} as a rational identity.
\end{proof}

\begin{proof}[Proof of Theorem~\ref{dhr:thm:gauss}]
On compact subsets of the stated domain,
$w_n=O((n+1)^{-1-\Rea d})$. For ordered nonnegative indices,
\[
 |\Delta(\bm n)|^2\leq\prod_{j=1}^r(n_j+1)^{2(j-1)}.
\]
Dropping the ordering gives a product of convergent majorant series
when $\Rea d>2r-2$. This proves absolute local normal convergence,
including that of moments of degrees at most $2r-2$.

Apply finite Cauchy--Binet to the two monic polynomial bases
$\fall{n}{i}$ and $(n+a)_j$, then pass to the normally convergent
limit. Both evaluation determinants equal the Vandermonde, so
\begin{equation}\label{dhr:eq:moment-det}
 Z_r=\det_{0\leq i,j<r}M_{ij},\qquad
 M_{ij}=\sum_{n\geq0}\fall{n}{i}(n+a)_j w_n.
\end{equation}
Shifting $n=i+k$ and applying the ordinary Gauss sum gives
\begin{equation}\label{dhr:eq:mixed-moment}
 M_{ij}=\frac{\Gamma(a+i+j)\Gamma(c+i)\Gamma(d-i-j)}
 {\Gamma(b-a-j)\Gamma(b-c)}.
\end{equation}
Indeed the hypergeometric numerator parameters are $a+i+j,c+i$,
the denominator parameter is $b+i$, and the excess is $d-i-j$.
Its real part is positive throughout the present index range.

Factor from row $i$ the quantity
\[
 R_i=\frac{\Gamma(a+i)\Gamma(c+i)\Gamma(d-i)}
 {\Gamma(b-a)\Gamma(b-c)},
\]
and from column $j$ the factor $(1-b+a)_j$. The remaining entry is
$(a+i)_j/(1-d+i)_j$. Lemma~\ref{dhr:lem:rational-det} applies with
$x_i=a+i$ and $t=1-b+c$. Since $x_k-x_i=k-i$, its Vandermonde
is $\prod_{j=0}^{r-1}j!$. The factorial identities
\[
 (1-b+a)_j=(-1)^j\frac{\Gamma(b-a)}{\Gamma(b-a-j)},\quad
 (1-b+c)_j=(-1)^j\frac{\Gamma(b-c)}{\Gamma(b-c-j)}
\]
and
\[
 (1-d+i)_{r-1}=(-1)^{r-1}
       \frac{\Gamma(d-i)}{\Gamma(d-i-r+1)}
\]
produce exactly \eqref{dhr:eq:gauss-product}; the total sign is one.
Perform the calculation first away from the auxiliary determinant's
poles, then extend by the identity theorem. The original sum and the
final product are holomorphic on the asserted domain.
\end{proof}

\begin{remark}[The one-variable boundary case]
For $r=1$, \eqref{dhr:eq:gauss-product} is simply
\[
 \sum_{n\geq0}w_n(a,c;b)
 =\frac{\Gamma(a)\Gamma(c)\Gamma(b-a-c)}
 {\Gamma(b-a)\Gamma(b-c)}.
\]
The subsequent $r=1$ resonance is the familiar zero-balanced Gauss
case. The continuation of interest is the uniform mechanism for all $r$,
not a claim of novelty for this boundary case.
\end{remark}

% END sections/02-determinant.tex

% BEGIN sections/03-resonance.tex
\section{First resonance and the correct finite part}
\label{dhr:sec:resonance}

Fix $r\geq1$, initially $a,c>0$, and set
\begin{equation}\label{dhr:eq:resonant-parameters}
 b_0=a+c+2r-2,\qquad b=b_0+u.
\end{equation}
Let $L(u)$ be any holomorphic function near zero with $L(0)=1$.
It is a normalization of \emph{each} particle's weight. Define
\begin{equation}\label{dhr:eq:amplitude-E}
 \begin{split}
 A(u)&=L(u)^r Z_{r-1}(a,c;b_0+u),\\
 E(u)&=\frac{\Gamma(r)\Gamma(a+r-1)\Gamma(c+r-1)\Gamma(1+u)^2}
 {\Gamma(r+u)\Gamma(a+r-1+u)\Gamma(c+r-1+u)}.
 \end{split}
\end{equation}
Both are holomorphic near zero, and $E(0)=1$. Although $A$ uses the
$(r-1)$-particle partition, its normalization is $L^r$, not $L^{r-1}$:
the escaping largest index contributes a normalization factor as well.

Write the contribution of a prescribed largest index as
\begin{equation}\label{dhr:eq:maximum-mass}
 \begin{split}
 K_n(u)={}&L(u)^r w_n(a,c;b_0+u)\\
 &\quad\times\sum_{0\leq m_1<\cdots<m_{r-1}<n}
 \Delta(\bm m)^2\prod_{i=1}^{r-1}
       (n-m_i)^2w_{m_i}(a,c;b_0+u).
 \end{split}
\end{equation}
For $r=1$ the inner sum is one, including at $n=0$; for $r>1$ it
vanishes if $n<r-1$. Each $K_n$ is holomorphic near zero.

\begin{theorem}[Largest-index subtraction]\label{dhr:thm:resonance}
For every fixed $\alpha$ with $\Rea\alpha>0$, the series
\begin{equation}\label{dhr:eq:D-def}
 D(u;\alpha)=\sum_{n\geq0}
 \left\{K_n(u)-A(u)(n+\alpha)^{-1-u}\right\}
\end{equation}
is locally normally convergent near $u=0$, and
\begin{equation}\label{dhr:eq:D-master}
 \boxed{D(u;\alpha)=A(u)
       \left\{\frac{E(u)}u-\zeta(1+u,\alpha)\right\}.}
\end{equation}
The apparent pole on the right is removable. Every fixed derivative
in $u$ can be taken termwise. The assertions hold jointly locally
holomorphically for $\Rea a,\Rea c,\Rea\alpha>0$.
\end{theorem}

\begin{proof}
Dividing the products in Theorem~\ref{dhr:thm:gauss} at particle
numbers $r$ and $r-1$, with the same $b=b_0+u$, gives
\begin{equation}\label{dhr:eq:meromorphic-master}
 L(u)^r Z_r(a,c;b_0+u)=\frac{A(u)E(u)}u.
\end{equation}
For instance, the ratio of the factors common to the two products
contains
\[
 \prod_{j=0}^{r-2}
 \frac{\Gamma(r-1+u-j)}{\Gamma(r+u-j)}
 =\frac{\Gamma(1+u)}{\Gamma(r+u)}.
\]
The remaining factor has $\Gamma(u)=\Gamma(1+u)/u$. These two
factors explain the square of $\Gamma(1+u)$ in $E$.
Equation~\eqref{dhr:eq:meromorphic-master} initially holds for
$\Rea u>0$ sufficiently small.

We next justify the subtraction independently of that product's
meromorphic continuation. Work on a compact parameter set and a small
closed disk $|u|\leq\delta<1/2$. The absolute $(r-1)$-particle sum
has a uniformly finite first moment:
\begin{equation}\label{dhr:eq:lower-first-moment}
 \sum_{m_1<\cdots<m_{r-1}}
 (1+m_1+\cdots+m_{r-1})
 |\Delta(\bm m)|^2\prod_i|w_{m_i}|<C.
\end{equation}
Indeed, the largest exponent in the separate-index majorant is
$-3-\Rea u$ before inserting a first moment, and at most
$-2+\delta<-1$ after insertion. This also proves the assertion
uniformly for complex $a,c$ in compact right-half-plane sets.

For $0\leq m_i<n$,
\[
 \left|1-\prod_i(1-m_i/n)^2\right|
 \leq\frac{2}{n}\sum_i m_i.
\]
The missing region $m_{r-1}\geq n$ is likewise bounded by the first
moment divided by $n$. It follows that the inner sum in
\eqref{dhr:eq:maximum-mass}, divided by $n^{2r-2}$, is
\[
 Z_{r-1}(a,c;b_0+u)+O(n^{-1})
\]
uniformly on that compact set. The assertion is immediate for $r=1$.
The uniform gamma-ratio expansion gives
\[
 w_n(a,c;b_0+u)=n^{-2r+1-u}\{1+O(n^{-1})\}.
\]
Consequently, for $n\geq1$,
\begin{equation}\label{dhr:eq:asymptotic-mass}
 K_n(u)=A(u)n^{-1-u}+O(n^{-2+\delta}).
\end{equation}
Replacing $n$ by $n+\alpha$ in the leading term changes only the
summable remainder, uniformly on compact $\alpha$-sets. The finitely
many initial terms are harmless. This proves local normal convergence
of \eqref{dhr:eq:D-def}. Cauchy's integral formula on a slightly
larger parameter disk proves normal convergence for every fixed
$u$-derivative; there is no unproved exchange of an infinite series
with differentiation at a pole.

For $\Rea u>0$, sum the leading counterterm by Hurwitz zeta and use
\eqref{dhr:eq:meromorphic-master}. This yields
\eqref{dhr:eq:D-master} on a nonempty open set. Both sides extend
holomorphically through zero, because $E(0)=1$ and the Hurwitz residue
is one. The identity theorem completes the proof, including joint
holomorphy in the auxiliary parameters.
\end{proof}

Put $A_j=A^{(j)}(0)$ and
\begin{equation}\label{dhr:eq:chi}
 \chi=E'(0)=2\psi(1)-\psi(r)-\psi(a+r-1)-\psi(c+r-1).
\end{equation}
In particular,
\begin{equation}\label{dhr:eq:residue}
 A_0=\prod_{j=0}^{r-2}
 \frac{j!\,\Gamma(a+j)\Gamma(c+j)\Gamma(r-j)}
 {\Gamma(c+2r-2-j)\Gamma(a+2r-2-j)}.
\end{equation}

\begin{corollary}[Hard-cutoff constant]\label{dhr:cor:cutoff}
For $a,c>0$,
\begin{equation}\label{dhr:eq:cutoff-constant}
 \boxed{\lim_{N\to\infty}\left[
 \sum_{0\leq n_1<\cdots<n_r\leq N}
   \Delta(\bm n)^2\prod_i w_{n_i}(a,c;b_0)
       -A_0\log N\right]=A_0\chi.}
\end{equation}
The corresponding spectral finite part is
\begin{equation}\label{dhr:eq:spectral-vs-cutoff}
 \FP_{u=0}\{L(u)^rZ_r(a,c;b_0+u)\}=A_1+A_0\chi.
\end{equation}
Thus the hard-cutoff constant is obtained from the spectral finite
part by subtracting $A_1$.
\end{corollary}
\begin{proof}
At zero, Theorem~\ref{dhr:thm:resonance} gives
\[
 D(0;\alpha)=A_0\{\chi-\gamma_0(\alpha)\}
            =A_0\{\chi+\psi(\alpha)\}.
\]
The partial sum of the subtracted counterterm is
$A_0\{\psi(N+1+\alpha)-\psi(\alpha)\}$.
Since $\psi(N+1+\alpha)=\log N+o(1)$, the auxiliary $\alpha$
cancels and proves \eqref{dhr:eq:cutoff-constant}.
Expanding $A(u)E(u)/u$ directly gives
\eqref{dhr:eq:spectral-vs-cutoff}.
\end{proof}

\begin{remark}[Why the normalization matters]
Changing $L$ while preserving $L(0)=1$ leaves the entire cutoff sum
unchanged. It changes $A_1$ by $rL'(0)A_0$ relative to the choice
$L\equiv1$, and changes the spectral finite part by the same amount.
Therefore identifying spectral and cutoff constants without specifying
the amplitude derivative cannot be a normalization-invariant statement.
This is a correction to that possible extrapolation, not an allegation
that the inspected manuscript makes the identification.
\end{remark}

% END sections/03-resonance.tex

% BEGIN sections/04-stieltjes.tex
\section{All-order Stieltjes identities and Hurwitz primitives}
\label{dhr:sec:stieltjes}

\subsection{The complete jet formula}
Let $e_j=E^{(j)}(0)$, so $e_0=1$ and $e_1=\chi$. All these
coefficients are explicit finite polygamma polynomials. More precisely,
put
\begin{equation}\label{dhr:eq:ell}
 \ell_j=2\psi^{(j-1)}(1)-\psi^{(j-1)}(r)
        -\psi^{(j-1)}(a+r-1)-\psi^{(j-1)}(c+r-1).
\end{equation}
Then $\log E(u)=\sum_{j\geq1}\ell_j u^j/j!$ and
$e_j=\Bell_j(\ell_1,\ldots,\ell_j)$. In particular,
\[
 e_2=\chi^2+\ell_2,\qquad
 e_3=\chi^3+3\chi\ell_2+\ell_3.
\]
The derivatives $A_j$ are obtained in exactly the same way from the
finite Gamma product \eqref{dhr:eq:amplitude-E}.

\begin{theorem}[Every fixed Stieltjes order]\label{dhr:thm:jets}
For every integer $k\geq0$, the following single-bracket series is
ordinary and absolutely convergent:
\begin{align}
 &\sum_{n\geq0}\left[
 K_n^{(k)}(0)-\frac1{n+\alpha}
 \sum_{j=0}^k\binom{k}{j}A_j[-\log(n+\alpha)]^{k-j}
 \right]\notag\\
 &\qquad=\sum_{j=0}^k\binom{k}{j}A_j
 \left\{\frac{e_{k-j+1}}{k-j+1}
       -(-1)^{k-j}\gamma_{k-j}(\alpha)\right\}.
 \label{dhr:eq:all-jets}
\end{align}
The result is locally holomorphic for the parameters in
Theorem~\ref{dhr:thm:resonance}. The constants $A_j,e_j$ are finite
Gamma/polygamma expressions; no infinite coefficient extraction is
needed to compute a fixed row.
\end{theorem}
\begin{proof}
Differentiate the already normally convergent identity
\eqref{dhr:eq:D-master}, rather than the divergent unsubtracted
partition. The expansion at zero is
\[
 \frac{E(u)}u-\zeta(1+u,\alpha)
 =\sum_{h\geq0}\left\{\frac{e_{h+1}}{h+1}
               -(-1)^h\gamma_h(\alpha)\right\}\frac{u^h}{h!}.
\]
Leibniz's rule gives the right side of \eqref{dhr:eq:all-jets}.
Differentiating $(n+\alpha)^{-1-u}$ gives the logarithmic polynomial
on the left. The normal-convergence statement proved in
Theorem~\ref{dhr:thm:resonance} supplies absolute convergence of the
resulting bracket and justifies every interchange.
\end{proof}

The first three values $D_k(\alpha)=\partial_u^kD(0;\alpha)$ are
\begin{align}
 D_0(\alpha)&=A_0\{\chi-\gamma_0(\alpha)\},\notag\\
 D_1(\alpha)&=A_1\{\chi-\gamma_0(\alpha)\}
                   +A_0\{e_2/2+\gamma_1(\alpha)\},\notag\\
 D_2(\alpha)&=A_2\{\chi-\gamma_0(\alpha)\}
          +2A_1\{e_2/2+\gamma_1(\alpha)\}
          +A_0\{e_3/3-\gamma_2(\alpha)\}.
 \label{dhr:eq:first-jets}
\end{align}
The factor $1/(h+1)$ in these rows comes from dividing $E(u)-1$ by
$u$ before differentiation. Replacing it by one would give an
incorrect higher-order table.

\subsection{Derivatives in the subtraction point}
The parameter $\alpha$ specifies the subtraction grid. It does not
move the particles or change their original weights. Differentiating
this auxiliary family nevertheless yields useful exact identities
between the corresponding subtracted sums and the usual Hurwitz
order ladder.

\begin{proposition}[Hurwitz derivative and shift laws]\label{dhr:prop:alpha}
For every integer $h\geq1$,
\begin{equation}\label{dhr:eq:alpha-derivative}
 \partial_\alpha^hD(u;\alpha)
 =(-1)^{h+1}A(u)(1+u)_h\zeta(1+u+h,\alpha).
\end{equation}
Moreover,
\begin{equation}\label{dhr:eq:alpha-shift}
 D(u;\alpha+1)-D(u;\alpha)=A(u)\alpha^{-1-u}.
\end{equation}
Every fixed $u$-derivative of these equations is valid, including
at $u=0$. In particular
$\partial_\alpha^hD_0(\alpha)=A_0\psi^{(h)}(\alpha)$.
\end{proposition}
\begin{proof}
Use the classical identities
$\partial_\alpha^h\zeta(s,\alpha)=(-1)^h(s)_h\zeta(s+h,\alpha)$
and $\zeta(s,\alpha+1)=\zeta(s,\alpha)-\alpha^{-s}$
\cite{dlmf:hurwitz}. Their substitution in
\eqref{dhr:eq:D-master} proves both formulas. Joint holomorphy, already
established, permits the indicated mixed derivatives. The last formula
also follows immediately from $D_0=A_0(\chi+\psi(\alpha))$.
\end{proof}

\subsection{Anchored antiderivatives and log-Gamma}
An anchored primitive fixes the integration constant, which is especially
important when differentiating in the spectral variable. Let
$\Rea\alpha_0,\Rea\alpha>0$, with the path in that half-plane.

\begin{theorem}[Every anchored primitive of the spectral jets]
\label{dhr:thm:primitive}
The identity
\begin{equation}\label{dhr:eq:primitive-master}
 \int_{\alpha_0}^{\alpha}D(u;t)\diff t
 =\frac{A(u)}u\left\{(\alpha-\alpha_0)E(u)
                  +\zeta(u,\alpha)-\zeta(u,\alpha_0)\right\}
\end{equation}
has a removable singularity at $u=0$. For every $k\geq0$,
\begin{align}
 \int_{\alpha_0}^{\alpha}D_k(t)\diff t
 =\sum_{j=0}^k\binom{k}{j}\frac{A_j}{k-j+1}
 \big[& (\alpha-\alpha_0)e_{k-j+1}\notag\\
 &+\zeta^{(k-j+1)}(0,\alpha)
 -\zeta^{(k-j+1)}(0,\alpha_0)\big].
 \label{dhr:eq:primitive-jets}
\end{align}
In particular,
\begin{equation}\label{dhr:eq:gamma-primitive}
 \int_{\alpha_0}^{\alpha}D_0(t)\diff t
 =A_0\left\{(\alpha-\alpha_0)\chi
       +\log\Gamma(\alpha)-\log\Gamma(\alpha_0)\right\}.
\end{equation}
\end{theorem}
\begin{proof}
For $u\ne0$, integrate the identity
$\partial_t\zeta(u,t)=-u\zeta(1+u,t)$.
This gives \eqref{dhr:eq:primitive-master}. The numerator vanishes
at zero because $E(0)=1$ and $\zeta(0,t)=\tfrac12-t$.
Its divided Taylor series has $h$th derivative at zero
\[
 \frac{(\alpha-\alpha_0)e_{h+1}
       +\zeta^{(h+1)}(0,\alpha)-\zeta^{(h+1)}(0,\alpha_0)}{h+1}.
\]
Leibniz's rule proves \eqref{dhr:eq:primitive-jets}. Finally,
$\zeta'(0,t)=\log\Gamma(t)-\tfrac12\log(2\pi)$
\cite{dlmf:hurwitz} gives \eqref{dhr:eq:gamma-primitive}.
\end{proof}

This theorem transports the classical Stieltjes primitive to a
many-variable resonant harmonic sum. It does not claim a new formula
for $\int\gamma_k(t)\diff t$ in isolation, nor does it integrate an
arbitrary many-variable weight with respect to a particle coordinate.
The distinction identifies precisely which antiderivative has been
established.

% END sections/04-stieltjes.tex

% BEGIN sections/05-harmonics.tex
\section{Integer and mixed harmonic specializations}
\label{dhr:sec:harmonic}

\subsection{A rational spectral partition at every particle number}
Take $a=c=1$, so $b_0=2r$, and choose
\begin{equation}\label{dhr:eq:integer-normalization}
 L(u)=\frac{\Gamma(2r+u)}{\Gamma(2r)}.
\end{equation}
At zero the particle weight is rational:
\[
 w_n(1,1;2r)=\frac1{(n+1)_{2r-1}}.
\]
Denote the normalized meromorphic partition by
$F_r(u)=L(u)^rZ_r(1,1;2r+u)$.

\begin{proposition}[Rationality of the normalized partition]
\label{dhr:prop:rational-partition}
For every $r\geq1$,
\begin{equation}\label{dhr:eq:rational-partition}
 \boxed{F_r(u)=
 \frac{\displaystyle\prod_{j=0}^{r-1}(j!)^3}{[(2r-1)!]^r}
 \prod_{k=0}^{r-1}\frac{(u+r+k)_{r-k}}{(u+k)_r}.}
\end{equation}
Its residue $A_0$ at zero is a positive rational number. The hard-cutoff
constant is
\begin{equation}\label{dhr:eq:integer-cutoff}
 A_0\{\gamma-3H_{r-1}\}.
\end{equation}
\end{proposition}
\begin{proof}
For $a=c=1$, the $j$th product factor in
\eqref{dhr:eq:gauss-product} is
\[
 (j!)^3\frac{\Gamma(r-1+u-j)}{\Gamma(2r-1+u-j)^2}.
\]
Pair one factor $\Gamma(2r+u)$ from $L^r$ with each of these,
use Gamma recurrences, and replace $j$ by $r-1-k$.
This gives \eqref{dhr:eq:rational-partition}. Finally,
$\chi=2\psi(1)-3\psi(r)=\gamma-3H_{r-1}$, and
Corollary~\ref{dhr:cor:cutoff} applies.
\end{proof}

The first partitions are
\begin{align}
 F_1(u)&=\frac{u+1}{u},\notag\\
 F_2(u)&=\frac{(u+3)^2}{36u(u+1)^2},\notag\\
 F_3(u)&=\frac{(u+4)(u+5)^3}
 {216000u(u+1)^2(u+2)^3(u+3)},\notag\\
 F_4(u)&=\frac{(u+6)^2(u+7)^4}
 {373403520000u(u+1)^2(u+2)^3(u+3)^4(u+4)^2}.
 \label{dhr:eq:small-partitions}
\end{align}
The corresponding residues are
\begin{equation}\label{dhr:eq:small-residues}
 A_0=1,\quad\frac14,\quad\frac1{10368},\quad
 \frac1{44789760000}\qquad(r=1,2,3,4).
\end{equation}
Although $F_r$ is rational, the individual normal-convergence amplitude
$A(u)$ is generally not rational. Its polygamma derivatives are exactly
what permit Stieltjes constants to appear in ordinary subtracted sums.

For calculation, the logarithmic derivatives are
\begin{align}
 (\log A)^{(k)}(0)
 &=r\psi^{(k-1)}(2r)
 +\sum_{j=0}^{r-2}\left\{\psi^{(k-1)}(r-j)
                    -2\psi^{(k-1)}(2r-1-j)\right\},
 \label{dhr:eq:integer-A-jets}\\
 \ell_k&=(-1)^k(k-1)!
       \{3H_{r-1}^{(k)}-\zeta(k)\},\qquad k\geq2.
 \label{dhr:eq:integer-E-jets}
\end{align}
Together with $\ell_1=\gamma-3H_{r-1}$, these are finite exact
coefficient recipes for every row of Theorem~\ref{dhr:thm:jets}.

\subsection{The derivatives as actual harmonic polynomials}
For a configuration $\bm n=(n_1,\ldots,n_r)$ define
\begin{equation}\label{dhr:eq:harmonic-log-jets}
 v_k(\bm n)=(-1)^k(k-1)!
  \sum_{i=1}^r\left\{H_{n_i+2r-1}^{(k)}-H_{2r-1}^{(k)}\right\}.
\end{equation}
Logarithmic differentiation of the normalized weight gives
\begin{equation}\label{dhr:eq:K-harmonic}
 K_n^{(k)}(0)=
 \sum_{0\leq n_1<\cdots<n_{r-1}<n_r=n}
 \frac{\Delta(\bm n)^2\,
  \Bell_k(v_1(\bm n),\ldots,v_k(\bm n))}
 {\displaystyle\prod_{i=1}^r(n_i+1)_{2r-1}}.
\end{equation}
This formula has no spectral derivatives left unevaluated. To see it,
for a single particle differentiate
$\log\Gamma(2r+u)-\log\Gamma(n+2r+u)$.
The polygamma shift formula gives the corresponding summand in
\eqref{dhr:eq:harmonic-log-jets}; the Bell identity
\eqref{dhr:eq:bell} then differentiates the product. Substitution in
\eqref{dhr:eq:all-jets} proves an explicit all-order family of
harmonic--logarithm--Stieltjes identities.

For a concrete first-order member, put
\begin{equation}\label{dhr:eq:W2}
 W_{m,n}=\frac{(n-m)^2}{(m+1)_3(n+1)_3},\qquad 0\leq m<n.
\end{equation}
At $r=2$, $A_0=1/4$, $A_1=(5/3-\gamma)/4$, and
$e_2=(\gamma-3)^2+3-\zeta(2)$. Thus
\begin{align}
 \sum_{n\geq0}\bigg[&\sum_{m=0}^{n-1}W_{m,n}
       (2H_3-H_{m+3}-H_{n+3})\notag\\
 &-\frac{5/3-\gamma-\log(n+1)}{4(n+1)}\bigg]
 =\frac14+\frac{\gamma^2-\zeta(2)}8+\frac{\gamma_1}{4}.
 \label{dhr:eq:explicit-stieltjes}
\end{align}
Here $\gamma_1=\gamma_1(1)$, and the empty inner sum at $n=0$ is zero.
The zero-order companion is
\begin{equation}\label{dhr:eq:explicit-stieltjes-zero}
 \sum_{n\geq0}\left\{\sum_{m<n}W_{m,n}-\frac1{4(n+1)}\right\}
 =-\frac34.
\end{equation}
These are ordinary convergent bracketed series, not a declaration of
finite-part integration in disguise.

The same example makes the normalization distinction numerical and exact:
\begin{equation}\label{dhr:eq:normalization-example}
 \FP_{u=0}F_2(u)=-\frac13,\qquad
 \lim_{N\to\infty}\left\{\sum_{m<n\leq N}W_{m,n}
                              -\frac14\log N\right\}
 =\frac{\gamma-3}{4}.
\end{equation}
Their difference is $A_1=(5/3-\gamma)/4$, as required. With $L=1$
the spectral finite part instead equals $(2\gamma-5)/4$; the hard
cutoff is unchanged.

\subsection{Ordinary all-order mixed harmonic sums}
A convergent, rather than resonant, specialization produces another
useful family. For an integer $B>2r$, define
\begin{equation}\label{dhr:eq:mixed-generator}
 \mathcal H_{r,B}(p,q)=
 \frac{Z_r(1+p,1+q;B)}{[\Gamma(1+p)\Gamma(1+q)]^r}.
\end{equation}
Near $p=q=0$, it is both a normally convergent sum and the explicit
Gamma product
\begin{align}
 \mathcal H_{r,B}(p,q)
 =\prod_{j=0}^{r-1}
 \frac{j!\,\Gamma(1+p+j)\Gamma(1+q+j)
              \Gamma(B-r-1-j-p-q)}
 {\Gamma(1+p)\Gamma(1+q)
              \Gamma(B-1-j-p)\Gamma(B-1-j-q)}.
 \label{dhr:eq:mixed-product}
\end{align}
The convergence region includes a neighborhood with
$\Rea(p+q)<B-2r$.

\begin{theorem}[Many-variable mixed harmonic identities]
\label{dhr:thm:mixed}
For a configuration $\bm n$, let
\[
 U_j(\bm n)=(-1)^{j-1}(j-1)!\sum_{i=1}^rH_{n_i}^{(j)},\qquad
 W_B(\bm n)=\frac{\Delta(\bm n)^2}{\prod_i(n_i+1)_{B-1}}.
\]
For all $k,l\geq0$,
\begin{align}
 &\sum_{n_1<\cdots<n_r}W_B(\bm n)
 \Bell_k(U_1,\ldots,U_k)\Bell_l(U_1,\ldots,U_l)\notag\\
 &\hspace{35mm}=\partial_p^k\partial_q^l
               \mathcal H_{r,B}(0,0).
 \label{dhr:eq:mixed-harmonic}
\end{align}
Every pure row ($k=0$ or $l=0$) is rational. Every mixed row lies
in $\Q[\zeta(2),\ldots,\zeta(k+l)]$; for zero or first total order
this means simply $\Q$.
\end{theorem}
\begin{proof}
Relative to its value at $p=q=0$, the summand is multiplied by
\[
 \prod_{i=1}^r\prod_{j=1}^{n_i}(1+p/j)(1+q/j).
\]
Its logarithmic derivatives in either variable are the $U_j$.
The Bell formula and local normal convergence prove
\eqref{dhr:eq:mixed-harmonic}. The pure generating function is
itself rational:
\begin{equation}\label{dhr:eq:pure-rational}
 \mathcal H_{r,B}(p,0)=
 \prod_{j=0}^{r-1}
 \frac{(j!)^2(1+p)_j}
 {\Gamma(B-1-j)(B-r-1-j-p)_r}.
\end{equation}
This follows by Gamma recurrences in \eqref{dhr:eq:mixed-product}.
It proves rationality of every pure derivative, not just the displayed
low orders. For derivatives involving both $p$ and $q$, only
$\log\Gamma(B-r-1-j-p-q)$ contributes to a mixed derivative
of the logarithm. Its derivatives of total order $h\geq2$ are
rational combinations of $\zeta(h)$ and finite generalized harmonic
numbers. Pure logarithmic derivatives are rational by
\eqref{dhr:eq:pure-rational}. Exponentiating with finite multivariate
Bell coefficients proves the asserted zeta algebra.
\end{proof}

For $r=2,B=6$, the product simplifies to
\begin{equation}\label{dhr:eq:H26}
 \mathcal H_{2,6}(p,q)=
 \frac{(1+p)(1+q)\Gamma(3-p-q)\Gamma(2-p-q)}
 {\Gamma(5-p)\Gamma(5-q)\Gamma(4-p)\Gamma(4-q)}.
\end{equation}
Set $V_{m,n}=(n-m)^2/[(m+1)_5(n+1)_5]$ and sum below over
$0\leq m<n$. Direct differentiation yields
\begin{align}
 \sum V_{m,n}&=\frac1{10368},\notag\\
 \sum V_{m,n}(H_m+H_n)&=\frac{29}{12\cdot10368},\notag\\
 \sum V_{m,n}\{(H_m+H_n)^2-H_m^{(2)}-H_n^{(2)}\}
 &=\frac{43}{8\cdot10368},\notag\\
 \sum V_{m,n}(H_m+H_n)^2
 &=\frac{517/144+2\zeta(2)}{10368},\notag\\
 \sum V_{m,n}\{H_m^{(2)}+H_n^{(2)}\}
 &=\frac{2\zeta(2)-257/144}{10368}.
 \label{dhr:eq:ordinary-examples}
\end{align}
For example, at zero the first logarithmic $p$-derivative is $29/12$,
the second pure logarithmic derivative is $-67/144$, and the mixed
logarithmic derivative is $2\zeta(2)-9/4$. These three numbers prove
all four nonconstant rows. They illustrate why a pure Bell combination
can be rational while the square of the ordinary harmonic sum is not.

% END sections/05-harmonics.tex

% BEGIN sections/06-depth-two.tex
\section{Rank-one harmonic dependence and depth-two closure}
\label{dhr:sec:depth}

\subsection{A theorem beyond consecutive factorial weights}
The next result is algebraic and does not require the Gamma product.
Let $w(x)\in\Q(x)$ have only simple poles, all at negative integers,
and suppose
\begin{equation}\label{dhr:eq:rational-weight-assumption}
 w(n)=\lambda n^{-(2r-1)}+O(n^{-2r}),\qquad
 \lambda\in\Q\setminus\{0\}.
\end{equation}
There is no positivity assumption. The pole locations need not be
consecutive. Define
\begin{equation}\label{dhr:eq:general-cutoff}
 S_{w,N}=\sum_{0\leq n_1<\cdots<n_r\leq N}
              \Delta(\bm n)^2\prod_iw(n_i),\qquad S_{w,-1}=0.
\end{equation}

\begin{theorem}[Rank-one harmonic dependence]\label{dhr:thm:rank-one}
There exist rational functions $P_w(N),Q_w(N)\in\Q(N)$ such that
for every integer $N\geq0$,
\begin{equation}\label{dhr:eq:rank-one-cutoff}
 \boxed{S_{w,N}=P_w(N)H_{N+1}+Q_w(N).}
\end{equation}
Both rational functions are bounded at infinity. At each finite pole,
$P_w$ has order at most $r-1$ and $Q_w$ has order at most $r$.
The possible poles are negative integers at most $-2$.
For $r=1$, $P_w=\lambda$.
\end{theorem}

\begin{proof}
For $0\leq k\leq2r-2$, partial fractions give
\begin{equation}\label{dhr:eq:moment-residues}
 n^kw(n)=\sum_{d\geq1}\frac{c_{k,d}}{n+d},\qquad
 \sum_dc_{k,d}=\lambda\ind_{k=2r-2}.
\end{equation}
Only finitely many coefficients are nonzero. There is no polynomial
part because $n^kw(n)$ is proper. The sum of residues is the coefficient
of $n^{-1}$ at infinity, proving the second equality.

Thus the finite moment $M_k(N)=\sum_{n=0}^Nn^kw(n)$ equals
\begin{align}
 M_k(N)&=R_k(N)+\lambda\ind_{k=2r-2}H_{N+1},\notag\\
 R_k(N)&=\sum_dc_{k,d}\left\{
       \sum_{\ell=2}^{d}\frac1{N+\ell}-H_{d-1}\right\}.
 \label{dhr:eq:finite-moments}
\end{align}
Every $R_k$ is a bounded rational function with simple poles only
at the indicated negative integers. Finite Cauchy--Binet gives
\begin{equation}\label{dhr:eq:hankel-cutoff}
 S_{w,N}=\det_{0\leq i,j<r}M_{i+j}(N).
\end{equation}
Among these entries, $i+j=2r-2$ occurs only at $i=j=r-1$.
Therefore
\begin{equation}\label{dhr:eq:rank-one-matrix}
 (M_{i+j})_{i,j=0}^{r-1}
 =(R_{i+j})_{i,j=0}^{r-1}
       +\lambda H_{N+1}\bm e_{r-1}\bm e_{r-1}^{\mathsf T}.
\end{equation}
The determinant is affine in the harmonic number, with
\begin{equation}\label{dhr:eq:PQ-det}
 P_w=\lambda\det_{0\leq i,j<r-1}R_{i+j},\qquad
 Q_w=\det_{0\leq i,j<r}R_{i+j}.
\end{equation}
The empty-determinant convention covers $r=1$. Products of $r-1$
or $r$ simple-pole entries prove the pole-order bounds, and boundedness
of every entry proves boundedness at infinity. This argument proves
\eqref{dhr:eq:rank-one-cutoff} at every cutoff, including $N<r-1$
where the determinant and configuration sum both vanish.
\end{proof}

\begin{remark}[The source of the depth collapse]
A generic determinant whose entries are affine in $H_{N+1}$ has
harmonic degree as large as $r$. That estimate is much too weak here.
The decay condition \eqref{dhr:eq:rational-weight-assumption} forces
the residue sum to vanish in every moment except the highest one.
Consequently the harmonic matrix has rank one \emph{and a single
nonzero entry}. This is the mechanism, not cancellation observed in
a few low-rank examples.
\end{remark}

\subsection{An explicit partial-fraction transform}
For later use put
\[
 T_{p,j}(z)=\sum_{m\geq1}\frac{z^m}{(m+j)^p},\qquad
 U_{p,j}(z)=\sum_{m\geq1}\frac{z^mH_m}{(m+j)^p},
 \qquad p\geq1,\ j\geq0.
\]
Shifting the summation index gives
\begin{equation}\label{dhr:eq:T-transform}
 T_{p,j}(z)=z^{-j}\left\{\Li_p(z)-
                    \sum_{k=1}^j\frac{z^k}{k^p}\right\}.
\end{equation}
For the harmonic sums the precise shifted formula is
\begin{align}
 U_{p,0}(z)&=\Li_{p,1}(z,1)+\Li_{p+1}(z),
 \label{dhr:eq:U-zero}\\
 U_{p,j}(z)&=z^{-j}\bigg\{\Li_{p,1}(z,1)
 -\sum_{k=1}^j\frac{z^kH_{k-1}}{k^p}
 -\sum_{\ell=1}^{j-1}\sum_{k=j+1}^{\infty}
             \frac{z^k}{k^p(k-\ell)}\bigg\},\quad j\geq1.
 \label{dhr:eq:U-shift}
\end{align}
For \eqref{dhr:eq:U-shift}, write $k=m+j$ and use
$H_m=H_{k-1}-\sum_{\ell=1}^{j-1}(k-\ell)^{-1}$.
The finite correction is $H_{k-1}$, not $H_k$.
Each last rational tail has a partial-fraction decomposition in
$k$ and $k-\ell$, hence involves only ordinary polylogarithms and
finite sums. Its starting index $j+1>\ell$ avoids the apparent pole.
Equations \eqref{dhr:eq:T-transform}--\eqref{dhr:eq:U-shift} are
explicit terminating reduction rules, not formal existence statements.

\begin{theorem}[All-particle depth-two closure]\label{dhr:thm:closure}
Under the assumptions of Theorem~\ref{dhr:thm:rank-one}, let
$K_{w,N}=S_{w,N}-S_{w,N-1}$ and
\begin{equation}\label{dhr:eq:general-generating}
 \mathcal G_w(z)=\sum_{N\geq0}K_{w,N}z^{N+1},\qquad |z|<1.
\end{equation}
Then $\mathcal G_w$ belongs to the $\Q(z)$-linear span of
\begin{equation}\label{dhr:eq:closure-space}
 1,\quad \Li_1(z),\ldots,\Li_r(z),\quad
 \Li_{1,1}(z,1),\ldots,\Li_{r-1,1}(z,1).
\end{equation}
For $r=1$ the depth-two part of this list is empty. Thus the depth
upper bound is two, independently of $r$, and the polylogarithmic
weight upper bound is $r$.
\end{theorem}

\begin{proof}
Set $m=N+1$. By \eqref{dhr:eq:rank-one-cutoff},
\begin{align}
 K_{w,N}={}&\{P_w(N)-P_w(N-1)\}H_{N+1}\notag\\
          &+Q_w(N)-Q_w(N-1)+\frac{P_w(N-1)}{N+1}.
 \label{dhr:eq:mass-difference}
\end{align}
The harmonic coefficient is a proper rational function of $m$,
with poles only at nonpositive integers and pole orders at most $r-1$.
It is zero for $r=1$. The nonharmonic coefficient is also proper,
with pole orders at most $r$: division by $m$ can increase the
order of $P_w(N-1)$ at zero only from at most $r-1$ to at most $r$.
All other poles already have the stated bound. The expressions at
$N=-1$ used in \eqref{dhr:eq:mass-difference} are well defined;
\eqref{dhr:eq:finite-moments} gives the empty moments there.

Decompose both coefficients into terms $(m+j)^{-p}$ with $j\geq0$.
For the nonharmonic part apply \eqref{dhr:eq:T-transform}, which
uses indices at most $r$. For the harmonic part apply
\eqref{dhr:eq:U-zero} and \eqref{dhr:eq:U-shift}. The resulting depth-two
index is at most $r-1$; any additional ordinary polylogarithm has
index at most $r$. This proves \eqref{dhr:eq:closure-space}.
All the series converge absolutely in the open disk, so the shifts
and finite partial fractions are legitimate.
\end{proof}

The equivalent identity
\begin{equation}\label{dhr:eq:Abel-cutoff}
 \mathcal G_w(z)=(1-z)\sum_{N\geq0}S_{w,N}z^{N+1}
\end{equation}
connects the cutoff and Abel descriptions. It is useful for checking
constants but is not needed to guess the depth bound.

\subsection{Derivatives, Euler primitives, and endpoint zeta reduction}
Let $\Theta=z\,d/dz$ and normalize $\Theta^{-1}$ by integration from
zero with the kernel $dt/t$. Then
\begin{equation}\label{dhr:eq:Euler-primitive}
 \Theta^{-q}\mathcal G_w(z)=
 \sum_{N\geq0}\frac{K_{w,N}}{(N+1)^q}z^{N+1},\qquad q\geq0.
\end{equation}
This also equals, for $q\geq1$,
\begin{equation}\label{dhr:eq:Euler-integral}
 \frac1{(q-1)!}\int_0^z\log^{q-1}(z/t)
                       \mathcal G_w(t)\frac{\mathrm dt}{t},
\end{equation}
initially along $0<t<z<1$, then as the analytic primitive of the
power series.

\begin{corollary}[Primitive closure and ordinary endpoint evaluations]
\label{dhr:cor:endpoint}
For every $q\geq0$, the function \eqref{dhr:eq:Euler-primitive} has
polylogarithmic depth at most two and weight at most $r+q$.
For $q\geq1$ its endpoint
\begin{equation}\label{dhr:eq:I-general}
 I_{w,q}=\sum_{N\geq0}\frac{K_{w,N}}{(N+1)^q}
\end{equation}
is an ordinary absolutely convergent sum and lies in the rational
linear span of
\[
 1,\qquad \zeta(k),\qquad \zeta(a)\zeta(b),
 \quad 2\leq k\leq r+q,\quad a,b\geq2,\quad a+b\leq r+q.
\]
No product of more than two zeta factors is needed for this endpoint
family. Ordinary derivatives and anchored ordinary antiderivatives
also retain depth at most two, although their rational prefactors
and exponent shifts change.
\end{corollary}

\begin{proof}
Divide the two proper rational coefficients in
\eqref{dhr:eq:mass-difference} by $m^q$. Their maximum pole orders
become at most $r-1+q$ in the harmonic part and $r+q$ in the other
part. The preceding proof therefore gives the asserted weight and
depth bounds.

Bounded rational functions have expansions at infinity, so
\eqref{dhr:eq:mass-difference} also gives
$K_{w,N}=P_w(\infty)/(N+1)+O(\log(N+2)/(N+1)^2)$.
This proves absolute convergence of \eqref{dhr:eq:I-general}.
The configuration sum itself is absolutely convergent after inserting
the maximum weight $(n_r+1)^{-q}$, by the same absolute
largest-index estimate used in Section~\ref{dhr:sec:resonance}.

The reduction rules give Laurent-polynomial, hence locally analytic
near $z=1$, coefficients of the polylogarithms. Only $\Li_1(z)$
and $\Li_{1,1}(z,1)=\tfrac12\log^2(1-z)$ diverge there. Convergence
forces their coefficients to vanish at one: a nonzero quadratic
logarithm cannot cancel a linear logarithm or a bounded term, and
a nonzero linear logarithm cannot cancel bounded terms. Those
vanishing-coefficient terms tend to zero. The other endpoint values
are $\zeta(p)$ and $\zeta(p,1)$. Lemma~\ref{dhr:lem:Euler} below
reduces the latter to the stated span.

For the last assertion, power-series differentiation multiplies
coefficients by affine factors in the index and shifts the exponent;
ordinary integration divides by such factors and shifts it oppositely.
A finite shift of $H_m$ changes it by a rational function. Polynomial
parts in the index, when they arise, sum to rational functions or
rational multiples of $\Li_1(z)$. Partial fractions of the remaining
proper parts again use \eqref{dhr:eq:T-transform}--\eqref{dhr:eq:U-shift},
with finite initial terms removed where needed. None of these operations
requires a third nested index.
\end{proof}

\begin{lemma}[Euler's height-one reduction]\label{dhr:lem:Euler}
For every integer $p\geq2$,
\begin{equation}\label{dhr:eq:Euler-zeta}
 \zeta(p,1)=\frac p2\zeta(p+1)
 -\frac12\sum_{j=1}^{p-2}\zeta(p-j)\zeta(j+1).
\end{equation}
\end{lemma}
\begin{proof}
We include a beta-integral proof of this classical identity so the
endpoint theorem has no unproved multiple-zeta reduction step.
For $|x|<1$, expand Euler's beta integral \cite{dlmf:beta}:
\begin{align*}
 B(1-x,y)-B(1,y)
 &=\int_0^1(t^{-x}-1)(1-t)^{y-1}\mathrm dt\\
 &=-\gamma-\psi(1-x)\\
 &\quad+\frac y2\big\{[\gamma+\psi(1-x)]^2
                          +\zeta(2)-\psi'(1-x)\big\}+O(y^2).
\end{align*}
The $y$-derivative at zero can be taken under the integral: the
factor $t^{-x}-1$ cancels its endpoint singularity at $t=1$.
Using $\log(1-t)/(1-t)=-\sum_{n\geq1}H_nt^n$ gives this derivative
as
\[
 -\sum_{n\geq1}H_n\left\{\frac1{n+1-x}-\frac1{n+1}\right\}
 =-\sum_{p\geq2}\zeta(p,1)x^{p-1}.
\]
The interchanges hold first for real $0<x<1$ by absolute convergence,
and then throughout the disk by holomorphy. Put
$f(x)=-\psi(1-x)-\gamma=\sum_{k\geq2}\zeta(k)x^{k-1}$.
Comparison with the beta expansion yields
\[
 \sum_{p\geq2}\zeta(p,1)x^{p-1}
 =\tfrac12\{f'(x)-\zeta(2)-f(x)^2\}.
\]
Equating coefficients of $x^{p-1}$ proves
\eqref{dhr:eq:Euler-zeta}.
\end{proof}

\subsection{A common translate and cyclotomic coordinates}
The rank-one mechanism survives translation of the entire simple-pole
grid. Define
\[
 H_n(a)=\sum_{j=0}^{n-1}\frac1{j+a},\qquad H_0(a)=0,\qquad a>0.
\]
Suppose all poles of $w(n)$ are simple and lie at $-a-d$ with
integers $d\geq0$, and the critical decay
\eqref{dhr:eq:rational-weight-assumption} holds. Use the monic moment
basis $(n+a)^k$. For each pole,
\[
 \sum_{n=0}^N\frac1{n+a+d}
 =H_{N+1}(a)+\sum_{\ell=1}^d\frac1{N+a+\ell}
                 -\sum_{\ell=0}^{d-1}\frac1{a+\ell}.
\]
The sum of residues still vanishes except in the top moment.
The proof of Theorem~\ref{dhr:thm:rank-one} therefore gives
\begin{equation}\label{dhr:eq:translated-cutoff}
 S_{w,N}=P_w(N,a)H_{N+1}(a)+Q_w(N,a),
\end{equation}
with the same pole-order bounds. For a rational weight depending
rationally on $a$, the coefficients are rational in $N,a$.

For rational $a=p/q\in(0,1]$ in lowest terms, put
$\omega=e^{2\pi i/q}$. The exact root-of-unity filters are
\begin{align}
 \sum_{n\geq0}\frac{z^{n+a}}{(n+a)^s}
 &=q^{s-1}\sum_{j=0}^{q-1}\omega^{-pj}
                       \Li_s(\omega^jz^{1/q}),
 \label{dhr:eq:single-filter}\\
 \sum_{n\geq0}\frac{z^{n+a}H_n(a)}{(n+a)^s}
 &=q^{s-1}\sum_{j,k=0}^{q-1}\omega^{-p(j+k)}
             \Li_{s,1}(\omega^jz^{1/q},\omega^k).
 \label{dhr:eq:double-filter}
\end{align}
These formulas hold for integers $s\geq1$ and $0<z<1$, and on any
common continuation with a fixed $q$th-root branch. To prove them,
write the top index as $qn+p$ and the inner index as $qm+p$.
The powers of $q$ from the denominators are $q^s$ and $q^{s+1}$;
the one and two residue filters divide by $q$ and $q^2$,
respectively. The strict inner inequality is preserved.

Finite index shifts and partial fractions, exactly as in
\eqref{dhr:eq:U-shift}, now prove depth at most two in level-$q$
colored polylogarithms for the generating function with exponent
$z^{N+a}$. The same conclusion holds for primitives that divide
its coefficients by $(N+a)^h$. No ordinary-zeta-only endpoint
reduction is asserted for general level $q$.

The shifted canonical weight $w(n)=1/(n+a)_{2r-1}$ is also covered
by the Gamma product: take the parameters there to be $(a,c)=(a,1)$
and $b_0=a+2r-1$. For $r=2$ its residue and cutoff constant are
\begin{equation}\label{dhr:eq:translated-cutoff-constant}
 A_0=\frac1{2a(a+1)},\qquad
 C_a=\frac{-2-\psi(a+1)}{2a(a+1)}.
\end{equation}
For example, at $a=1/2$,
\begin{equation}\label{dhr:eq:half-cutoff}
 \lim_{N\to\infty}\left[
 \sum_{0\leq m<n\leq N}
 \frac{(n-m)^2}{(m+\tfrac12)_3(n+\tfrac12)_3}
       -\frac23\log N\right]
 =\frac23\{\gamma+2\log2-4\}.
\end{equation}
This gives a direct digamma/cyclotomic bridge without changing the
finite-cutoff depth mechanism.

% END sections/06-depth-two.tex

% BEGIN sections/07-explicit.tex
\section{Explicit polylogarithms and ordinary multiple sums}
\label{dhr:sec:explicit}

In this section use the canonical critical weight
$w_r(n)=1/(n+1)_{2r-1}$. Abbreviate $S_{r,N}=S_{w_r,N}$,
$K_{r,N}=S_{r,N}-S_{r,N-1}$, and
\begin{equation}\label{dhr:eq:canonical-GI}
 G_r(z)=\sum_{N\geq0}K_{r,N}z^{N+1},\qquad
 I_{r,q}=\sum_{0\leq n_1<\cdots<n_r}
 \frac{\Delta(\bm n)^2}
 {(n_r+1)^q\prod_i(n_i+1)_{2r-1}},\quad q\geq1.
\end{equation}
Thus $I_{r,q}=\Theta^{-q}G_r(1)$, an ordinary convergent endpoint.
The one-variable case is $G_1=\Li_1$ and $I_{1,q}=\zeta(q+1)$.

\subsection{Two variables: complete cutoff and generating function}
The three moments needed for $r=2$ are
\begin{align}
 M_0(N)&=\frac{(N+1)(N+4)}{4(N+2)(N+3)},\notag\\
 M_1(N)&=\frac{N(N+1)}{4(N+2)(N+3)},\notag\\
 M_2(N)&=H_{N+1}-\frac{11}{4}
                      +\frac1{2(N+2)}+\frac9{2(N+3)}.
 \label{dhr:eq:r2-moments}
\end{align}
They follow directly from the residues of $n^k/(n+1)_3$.
Taking $M_0M_2-M_1^2$ gives the finite identity
\begin{equation}\label{dhr:eq:r2-cutoff}
 \boxed{S_{2,N}=
 \frac{(N+1)(N+4)}{4(N+2)(N+3)}H_{N+1}
 -\frac{(N+1)^2(3N+8)}{4(N+2)^2(N+3)}.}
\end{equation}
This proves the cutoff formula without any asymptotic summation.
It also independently recovers the residue $1/4$ and constant
$(\gamma-3)/4$.

Applying the exact transforms in Section~\ref{dhr:sec:depth} yields
\begin{align}
 G_2(z)={}&\frac{(1-z)^2}{4z^2}\log^2(1-z)
        +\frac{z-1}{2z}\Li_2(z)\notag\\
 &-\frac{6-5z}{4z^2}\log(1-z)+\frac34-\frac3{2z}.
 \label{dhr:eq:r2-G}
\end{align}
The expression is analytic at zero and starts $z^2/144$,
corresponding to the unique first configuration $(0,1)$.
Every negative-power coefficient cancels exactly.
At $z=1/2$, the classical dilogarithm reflection formula
\cite{dlmf:polylog} gives
\begin{equation}\label{dhr:eq:G-half}
 \sum_{0\leq m<n}\frac{(n-m)^2\,2^{-(n+1)}}{(m+1)_3(n+1)_3}
 =\frac12\log^22+\frac72\log2-\frac{\pi^2}{24}-\frac94.
\end{equation}
At the other endpoint, \eqref{dhr:eq:r2-G} gives
\begin{equation}\label{dhr:eq:Abel-constant}
 G_2(z)=-\frac14\log(1-z)-\frac34+o(1),\qquad z\uparrow1.
\end{equation}
The Abel constant is $-3/4$, not the hard-cutoff constant
$(\gamma-3)/4$. This is consistent with
$\sum_{N\geq0}H_{N+1}z^{N+1}=-\log(1-z)/(1-z)$.

\begin{theorem}[An ordinary double-sum evaluation]
\label{dhr:thm:first-evaluation}
\begin{equation}\label{dhr:eq:I21}
 \boxed{\displaystyle
 \sum_{0\leq m<n}
 \frac{(n-m)^2}{(m+1)_3(n+1)_3(n+1)}
 =\zeta(3)-\frac{\pi^2}{24}-\frac34.}
\end{equation}
\end{theorem}
\begin{proof}
We give a telescoping proof in addition to the polylogarithmic
primitive method. Put $v=N+1$. Differencing \eqref{dhr:eq:r2-cutoff}
gives
\[
 K_{2,v-1}=\frac{4(v+1)H_v+v^3-3v^2-6v}
 {4v(v+1)^2(v+2)}.
\]
Partial fractions of $K_{2,v-1}/v$ give a harmonic coefficient
\[
 -\frac3{4v}+\frac1{v+1}-\frac1{4(v+2)}+\frac1{2v^2}
\]
and a nonharmonic coefficient
\[
 -\frac3{4v}+\frac5{4(v+1)}-\frac1{2(v+2)}
                                      +\frac1{2(v+1)^2}.
\]
Group the first three terms of each coefficient into telescoping
differences before summing. Absolute convergence then permits
\begin{align*}
 \sum_{v\geq1}\frac{H_v}{v(v+1)}&=\zeta(2),&
 \sum_{v\geq1}\frac{H_v}{(v+1)(v+2)}&=1,&
 \sum_{v\geq1}\frac{H_v}{v^2}&=2\zeta(3).
\end{align*}
The first two identities follow by summing over the inner index in
$H_v$ and telescoping; the third is
$\zeta(2,1)+\zeta(3)=2\zeta(3)$ by
Lemma~\ref{dhr:lem:Euler}. The harmonic part therefore sums to
$-3\zeta(2)/4+1/4+\zeta(3)$. The nonharmonic part sums to
$\zeta(2)/2-1$. Their total is
$\zeta(3)-\zeta(2)/4-3/4$, proving the theorem.
\end{proof}

Two further primitives give
\begin{align}
 I_{2,2}&=\frac74-\frac{\pi^2}{16}-\frac32\zeta(3)
                            +\frac{\pi^4}{144},\label{dhr:eq:I22}\\
 I_{2,3}&=-\frac52+\frac{17\pi^2}{96}+\zeta(3)
             -\frac{\pi^4}{96}-\frac{\pi^2\zeta(3)}{12}
             +\frac32\zeta(5).
 \label{dhr:eq:I23}
\end{align}
These follow by dividing the same rational-plus-harmonic mass by
$v^2$ and $v^3$, applying \eqref{dhr:eq:T-transform}--\eqref{dhr:eq:U-shift},
and then \eqref{dhr:eq:Euler-zeta}. The delivered exact formula file
contains the full primitive reductions, not just their endpoints.

\subsection{Three variables: a depth-two formula of weight three}
For $r=3$, the determinant gives
\begin{equation}\label{dhr:eq:r3-P}
 P_3(N)=\frac{N(N+1)^2(N^2+17N+54)}
 {10368(N+3)^2(N+4)^2(N+5)}
\end{equation}
and
\begin{equation}\label{dhr:eq:r3-Q}
 Q_3(N)=-\frac{N(N+1)^3(9N^3+125N^2+514N+648)}
 {20736(N+2)(N+3)^3(N+4)^2(N+5)},
\end{equation}
with $S_{3,N}=P_3(N)H_{N+1}+Q_3(N)$. Hence its hard-cutoff
constant is $(\gamma-9/2)/10368$, again agreeing with the independent
Gamma resonance calculation.

For compact notation let $L_j=\Li_j(z)$ and
$D_j=\Li_{j,1}(z,1)$. The generating function is
\begin{align}
 G_3(z)={}&-\frac{5(z-1)^2(7z-4)}{1728z^4}D_1
        +\frac{(z-1)(z+1)}{144z^3}D_2\notag\\
 &-\frac{117z^3+95z^2-513z+300}{10368z^4}L_1\notag\\
 &-\frac{5(z-1)(7z-18)}{1728z^3}L_2
        +\frac{z-1}{144z^2}L_3
        -\frac{281z^2-78z-200}{6912z^3}.
 \label{dhr:eq:r3-G}
\end{align}
Its first nonzero coefficient is $z^3/54432000$, the configuration
$(0,1,2)$. Thus three ordered summation variables have not produced
depth three. The theorem explains this before any coefficient
calculation is performed.

The first three ordinary primitive endpoints are
\begin{align}
 I_{3,1}&=\frac{1903}{124416}
          -\frac{35\pi^2}{62208}-\frac7{864}\zeta(3),
 \label{dhr:eq:I31}\\
 I_{3,2}&=-\frac{6775}{373248}
          +\frac{67\pi^2}{186624}+\frac7{576}\zeta(3),
 \label{dhr:eq:I32}\\
 I_{3,3}&=\frac{100297}{4478976}
          -\frac{403\pi^2}{1492992}-\frac{131}{7776}\zeta(3)
          +\frac{\pi^4}{186624}.
 \label{dhr:eq:I33}
\end{align}
In particular, \eqref{dhr:eq:I31} is the explicit triple sum
\[
 \sum_{0\leq i<j<k}
 \frac{(j-i)^2(k-i)^2(k-j)^2}
 {(i+1)_5(j+1)_5(k+1)_5(k+1)}.
\]
The displayed value is a reduction of this entire sum, not a bound or
an asymptotic approximation.

\subsection{Four variables and the uniform algorithm}
A first four-variable endpoint is
\begin{equation}\label{dhr:eq:I41}
 I_{4,1}=\frac{630863}{268738560000}
 +\frac{14717\pi^2}{67184640000}
 -\frac{487}{124416000}\zeta(3)
 +\frac{\pi^4}{497664000}.
\end{equation}
The definition \eqref{dhr:eq:canonical-GI} specifies its six squared
difference factors and four degree-seven factorial denominators.
The numbers in these expressions can cancel substantially; direct
low-precision evaluation is therefore a poor way to assess them.
The exact rational reduction precedes the high-precision diagnostic.

For any requested $r$ and primitive order $q$, there is a terminating
algorithm: compute the $2r-1$ finite moments by residues, take the two
determinants \eqref{dhr:eq:PQ-det}, difference the cutoff, divide by
$(N+1)^q$, and apply the finite partial-fraction transform. It outputs
both a functional polylogarithmic identity and, when $q\geq1$, its
ordinary zeta endpoint. The implementation in \texttt{code/reductions.py}
uses exact rational arithmetic. It is not restricted in its mathematics
to the ranks and primitive orders printed here, although symbolic
expression growth limits convenient computation at large rank.

The fact that some displayed endpoints omit the highest permitted
weight is an additional cancellation in those cases. We assert the
proved upper bound $r+q$, not a universal sharp-weight formula and
not independence of the surviving zeta coordinates.

% END sections/07-explicit.tex

% BEGIN sections/08-audit.tex
\section{Source audit, corrections, and integration status}
\label{dhr:sec:audit}

\subsection{What was actually inspected}
The mathematical source anchor for the selected chapter reads is the
\repo\ commit
\begin{center}
\texttt{87c7db629552afbc2e803001526debe3430a31b0}.
\end{center}
The inspected material comprises the canonical manuscript README,
the mixed-Dougall chapter (including its full ending), the initial
240 lines of the cyclotomic-quotient chapter, the incoming-directory
inventory, and the first 180 lines of the incoming intake procedure.
Targeted repository searches supplied additional context about existing
Gauss, Stieltjes, tail, and Selberg-related work. Search snippets were
not treated as complete proofs of the unseen material.

The repository changed during retrieval: later search results carried
a different commit. This report therefore records a source anchor,
not an assertion that the whole repository was frozen throughout the
session. Six ZIP archives were visible in the incoming inventory at
the time it was read. Their binary contents were not unpacked or
analytically audited in this work. File names, including date strings,
do not establish their mathematical content. No statement here should
be read as a certification of those archives or of the entire four-volume
manuscript.

No substantive false theorem was found in the fully inspected
mixed-Dougall proof. Its all-order mixed derivatives and tail identities
remain prior repository material, not contributions claimed by this
report. The already corrected Choi normalization is likewise not a
newly discovered correction here \cite{repo:readme,repo:dougall}.

\subsection{Specific correction and normalization safeguards}
The principal correction supplied by the present analysis is the
\emph{transfer rule} between two possible definitions of a finite
constant. For this family the correct formula is
\[
 \text{hard-cutoff constant}
 =\text{spectral finite part}-A'(0).
\]
Equation~\eqref{dhr:eq:normalization-example} is an explicit
counterexample to omitting the last term. The rule corrects that
possible extrapolation; it is not attributed to any particular source
statement that was not observed.

Two further safeguards are necessary when implementing the formulas.
The leading amplitude uses $L(u)^r$, even though its partition has
$r-1$ particles. In the shifted harmonic transform
\eqref{dhr:eq:U-shift}, the finite correction is $H_{k-1}$.
Both choices follow from direct derivations, and the delivered
coefficient checks are sensitive to the shift convention. The
ordered-configuration normalization also precludes an extra factor
of $r!$.

The finite cutoff, the Abel constant, the spectral finite part, and the
value of an ordinary subtracted bracket are four distinct normalizations.
For the canonical two-particle example they are, respectively,
\[
 \frac{\gamma-3}{4},\qquad -\frac34,\qquad
 -\frac13\ \text{for the normalization \eqref{dhr:eq:integer-normalization}},
 \qquad -\frac34\ \text{at subtraction point }\alpha=1.
\]
The numerical equality of the second and fourth constants in this
example does not identify all four operations.

\subsection{Claim ledger}
\begin{center}
\small
\begin{tabular}{>{\raggedright\arraybackslash}p{0.28\textwidth}>{\raggedright\arraybackslash}p{0.30\textwidth}>{\raggedright\arraybackslash}p{0.31\textwidth}}
\toprule
Result & Status in this report & Scope qualification\\
\midrule
Gauss--Vandermonde product & Complete determinant proof & Classical inputs; no priority claim\\
Largest-index finite part & Complete analytic proof & First logarithmic resonance only\\
All Stieltjes orders & Complete normal-convergence proof & Fixed finite derivatives; full brackets\\
Harmonic mixed generator & Complete Gamma-product proof & Integer $B>2r$, local parameter germ\\
Rank-one harmonic cutoff & Complete algebraic proof & Simple negative-integer poles and critical decay\\
Depth-two closure & Complete terminating reduction & Upper bound, not minimal depth\\
Zeta endpoint evaluations & Complete reduction with Euler lemma & Ordinary positive primitive orders\\
Common translated grid & Complete residue/filter proof & Common grid, not arbitrary unrelated shifts\\
$S_6$, $S_8$ & Not resolved & Prior conjectural status unchanged\\
Proof-assistant status & Not formalized & No Lean or other kernel certificate\\
\bottomrule
\end{tabular}
\end{center}

The source's distinction between formal relation systems and actual
period independence is preserved \cite{repo:quotients}. A finite
rank calculation or a failed integer-relation search would not prove
that an endpoint lies outside a smaller period space. No such
nonreduction assertion is used here.

\subsection{Proposed integration}
A suitable standalone destination is
\begin{center}
\small\path{Analysis/Polylogarithms/docs/reports/determinantal-harmonic-resonance/}.
\end{center}
Its thematic cross-links are the Gamma--harmonic resonance material,
Hurwitz--Stieltjes calculus, and depth/cyclotomic chapters. The report
contains a complete article and can be reviewed independently before
any canonical chapter integration. All theorem and equation labels
start with \texttt{dhr:} to reduce collisions.

The incoming procedure explicitly separates preservation, replay,
analytic review, and canonical integration \cite{repo:intake}. This
delivery does not modify the GitHub repository, apply a patch, or
promote its own results into the canonical manuscript. The separate
\texttt{integration/INTEGRATION.md} records a suggested reading order
and chapter mapping. The delivered source files, not this paragraph,
are the proposed integration artifacts.

% END sections/08-audit.tex

% BEGIN sections/09-research.tex
\section{Further research questions}
\label{dhr:sec:research}

\subsection{Higher resonances: several divergent moment entries}
The first resonance $b-a-c=2r-2$ is special because only one particle
can escape at logarithmic order while the lower ensemble retains a
finite first moment. At $b-a-c=2r-2-M$ with $M\geq1$, a larger
block of moments is divergent and the Gamma product can have higher
pole order. Determine a canonical largest-index or multi-index
subtraction, including all overlap terms, and compare its hard-cutoff
polynomial in $\log N$ with its spectral Laurent coefficients.
A successful extension must prove normal convergence after the full
subtraction, not merely expand the Gamma product.

A precise initial target is the next resonance for $r=2$, where one
can explicitly compare the two divergent moments in the cutoff
Hankel determinant with the two relevant gamma poles. This should
clarify which derivatives of residue amplitudes enter beyond the
first-order rule \eqref{dhr:eq:spectral-vs-cutoff}.

\subsection{Repeated rational poles and the rank of harmonic data}
The simple-pole hypothesis in Theorem~\ref{dhr:thm:rank-one} is
essential to the proof. With repeated poles, moments contain
$H_{N+1}^{(2)},H_{N+1}^{(3)},\ldots$ as well as $H_{N+1}$.
Describe the matrix subspace spanned by these harmonic coefficients.
Can its rank and support give a sharp \emph{formal} upper bound on
the polylogarithmic depth of the cutoff generator?

One concrete problem is to classify all rational weights for which
this harmonic matrix pencil has rank one despite repeated poles.
Such a classification could extend the current depth-two theorem
without relying on generic determinant-degree estimates.

\subsection{Closed all-rank coefficients and discrete orthogonal polynomials}
The determinant formulas for $P_r(N)$ and $Q_r(N)$ are exact at every
rank, but expression growth is substantial. Find product formulas,
short recurrences, or a discrete-orthogonal-polynomial interpretation
for the rational coefficients in the canonical case. The factors of
$P_r$ displayed at ranks two through four suggest more structure
than the determinant proof alone explains.

A useful target is a recurrence in $r$ with rational coefficients in
$N$ that generates both $P_r$ and $Q_r$ and preserves the initial
vanishing $S_{r,N}=0$ for $N<r-1$. A companion recurrence for
$I_{r,q}$ could explain systematic weight drops in
\eqref{dhr:eq:I31}--\eqref{dhr:eq:I41}.

\subsection{Mixed spectral and shape derivatives}
Theorem~\ref{dhr:thm:resonance} gives joint holomorphy in $a,c,u$.
One can therefore differentiate in $a,c$ as well as $u$, provided
one keeps track of the moving relation $b_0=a+c+2r-2$.
Develop explicit harmonic/polygamma formulas for these tangential
shape derivatives and compare them with derivatives taken at fixed
$b$. The distinction changes the counterterm and is likely important
for integration with the existing mixed-Dougall calculus.

A targeted problem is to derive a compact three-parameter Bell table
whose rows express mixed generalized harmonic sums in terms of
$\gamma_k(\alpha)$ and polygamma values at $a+r-1,c+r-1$.
The analytic permission to differentiate is proved here; the most
useful arithmetic basis for the resulting rows is not selected.

\subsection{Rational grids beyond a common translate}
The root filters in \eqref{dhr:eq:single-filter} and
\eqref{dhr:eq:double-filter} handle one common residue class.
For rational poles from several classes modulo a common denominator,
the moment matrix can involve several independent shifted harmonic
families. Determine whether its harmonic rank, rather than the number
of summation variables, controls colored depth.

The first test case should use two residue classes with separate
harmonic coefficients and an explicitly computed matrix pencil.
Any resulting endpoint reduction must state its cyclotomic level and
its precise relation system; ordinary-zeta closure should not be
assumed from the common-integer-grid case.

\subsection{Changing the Vandermonde exponent}
The square $\Delta^2$ is what permits Cauchy--Binet and a moment
determinant. For $\Delta^{2\kappa}$ with integer $\kappa>1$, one
can still expand finitely in monomials, but the single-entry argument
no longer applies automatically. Determine the first resonant decay,
its subtraction amplitude, and a useful depth bound. A comparison
with known discrete Selberg-type sums should precede any novelty
claim in that direction.

\subsection{Functional certificates and arithmetic conjectures}
The current identities live in a tightly specified depth-two family.
It is reasonable to ask whether translated or colored variants supply
new functional relations useful for the manuscript's arithmetic
conjectures. There is no result here showing that they settle $S_6$
or $S_8$. A productive target would be an exact functional certificate
whose evaluation gives a previously missing relation, not another
high-precision residual or an unsupported claim of period dependence.

\subsection{Formalization and literature reconciliation}
The rational determinant lemma, finite Cauchy--Binet identity, residue
sum criterion, and terminating transform are natural first targets
for formalization. The analytic normal-convergence argument should
be formalized separately, with its compact-uniform first-moment
estimate and Cauchy derivative transfer made explicit.

In parallel, a systematic literature audit should identify which
forms of the Gauss--Vandermonde product and height-one generating
functions already occur in discrete determinant and Euler-sum
literature, then isolate the exact scope of the resonant cutoff and
rank-one harmonic statements. The present report supplies proofs and
reproducible identities; it does not turn an incomplete priority
search into a novelty theorem.

% END sections/09-research.tex

\appendix
% BEGIN sections/10-replay.tex
\section{Reproducible computation and proof boundaries}
\label{dhr:sec:replay}

\subsection{What is in the package}
The article consists of \texttt{article.tex}, the files under
\texttt{sections/}, and \texttt{references.tex}. The standalone
\texttt{article.pdf} is the compiled reading artifact.
The mathematical software is
\begin{center}
\begin{tabular}{lp{0.60\textwidth}}
\toprule
File & Purpose\\
\midrule
\texttt{code/reductions.py} & Exact canonical cutoff, polylogarithm, and endpoint reducer\\
\texttt{code/verify\_exact.py} & Finite rational/symbolic checks and formula export\\
\texttt{code/verify\_numeric.py} & Separate floating-point diagnostics\\
\bottomrule
\end{tabular}
\end{center}
The dependencies used in the recorded run are SymPy 1.14.0 and mpmath
1.3.0. Python 3.10 or later is suitable for the syntax in the scripts;
the recorded run used Python 3.13. The scripts write their own output
records, so replay should be performed on a working copy when the
original delivery records need to be preserved.

Run from the package root:
\begin{verbatim}
python -m pip install -r requirements.txt
python code/verify_exact.py
python code/verify_numeric.py
lualatex -interaction=nonstopmode -halt-on-error article.tex
lualatex -interaction=nonstopmode -halt-on-error article.tex
\end{verbatim}
A third LaTeX pass is harmless if a local build reports changed labels.
No network access is needed by the mathematical scripts after the
Python dependencies have been installed.

\subsection{Exact assertions}
The recorded exact suite passes 732 assertions. Grouping closely
related checks gives the following accounting:
\begin{center}
\small
\begin{tabular}{lr}
\toprule
Check family & Assertions\\
\midrule
Rational factorial determinant & 18\\
Gamma product versus mixed moment determinant & 16\\
Finite ordered sums, Cauchy--Binet, and finite moments & 216\\
Rational partitions, residues, and cutoff constants & 24\\
Polylogarithmic coefficients and removable origin terms & 416\\
Nonconsecutive-pole residue and moment checks & 32\\
Printed endpoint values and $G_2$ & 4\\
Intentionally corrupted functional formula rejected & 1\\
Mixed harmonic coefficient rows & 4\\
Printed Stieltjes row algebra & 1\\
\midrule
Total & 732\\
\bottomrule
\end{tabular}
\end{center}
These are finite exact checks. They detect arithmetic, index, and
normalization mistakes in explicit cases; they do not replace the
all-index proofs of Sections~\ref{dhr:sec:determinant}--\ref{dhr:sec:depth}.
The corruption control changes $G_2$ by $z^2$ and verifies that the
identity check rejects it.

The file \texttt{results/exact\_formulas.json} contains ranks
$r=1,2,3,4$, their cutoff coefficients, rational spectral partitions,
residues, generating functions of primitive orders $q=0,1,2,3$, and
endpoint values for $q=1,2,3$. In that file \texttt{L1}, \texttt{L2},
\ldots denote $\Li_1(z),\Li_2(z),\ldots$; \texttt{D1}, \texttt{D2},
\ldots denote $\Li_{1,1}(z,1),\Li_{2,1}(z,1),\ldots$.
The coefficient expressions use exact rational numbers and SymPy
syntax. They are not sampled numerical fits.

\subsection{Numerical diagnostics}
The recorded suite has 85 diagnostics, all at 100 decimal digits of
\emph{working precision}:
\begin{center}
\small
\begin{tabular}{lr}
\toprule
Diagnostic family & Cases\\
\midrule
Generating functions and Euler primitives from direct lower moments & 24\\
Ordinary Gauss product, nonintegral $a,c$ & 8\\
Convergent endpoints with an asymptotic tail & 12\\
Spectral jets, $r=1,\ldots,4$, orders $0,\ldots,4$, two $\alpha$ values & 40\\
Halfpoint logarithm/dilogarithm value & 1\\
\midrule
Total & 85\\
\bottomrule
\end{tabular}
\end{center}
The largest recorded scaled discrepancy is approximately
$2.75\times10^{-32}$. Individual tolerances and absolute as well as
scaled discrepancies are stored in \texttt{results/numerical\_checks.json}.
It would be incorrect to describe all checks as agreements to 100
digits: working precision is not an error bound.

For disk evaluations, the left side is formed from directly accumulated
lower moments and the right side from convergent polylogarithm series.
For ordinary Gamma products, 1500 terms are used at a large positive
excess. Endpoint diagnostics combine a 400-term direct head with a
30th-order Euler--Maclaurin expansion of the harmonic/rational tail,
summed by Hurwitz-zeta values and order derivatives. These tail
approximations are diagnostic, not interval-certified enclosures.

The spectral-jet left side uses an independent finite telescoping
formula instead of the infinite Gamma product. Namely,
\begin{align}
 &\sum_{n=0}^N(n+1)_j
       \frac{\Gamma(n+1)}{\Gamma(n+2r+u)}\notag\\
 &\qquad=
 \frac{\Gamma(j+1)/\Gamma(2r-1+u)
             -\Gamma(N+j+2)/\Gamma(N+2r+u)}{2r-j-2+u}.
 \label{dhr:eq:finite-jet-telescope}
\end{align}
To prove it, subtract successive terms of
$\Gamma(n+j+1)/\Gamma(n+2r-1+u)$; their difference is
$(2r-j-2+u)\Gamma(n+j+1)/\Gamma(n+2r+u)$.
At $j=2r-2,u=0$ the quotient has a removable singularity, which
is evaluated by Taylor coefficients rather than by division of two
rounded near-zero numbers. Transforming powers of $n$ to rising
factorials then gives the full finite moment determinant and its
spectral derivatives. The diagnostic cutoff is $N=10^{40}$;
it is used through this finite formula, not by enumerating that
many summands.

The translated root-of-unity extension is proved algebraically in the
article; the recorded 85-case numerical suite does not separately test
that extension. This is part of the verification scope, not a gap
filled by an unstated computation.

\subsection{Proof, calculation, and compilation are different evidence}
A clean LaTeX build establishes that the document compiles. Finite
symbolic assertions establish the corresponding finite identities.
Floating-point comparisons test implementations and motivate scrutiny.
The all-order analytic and algebraic conclusions are justified by the
proofs in the article. None of the three forms of computational evidence
alone confers theorem-prover status or proves global novelty.

% END sections/10-replay.tex

% BEGIN references.tex
\begin{thebibliography}{99}
\addcontentsline{toc}{section}{References}

\bibitem{repo:readme}
ProveIt Contributors,
\emph{Polylogarithms and their Arithmetic Bridges}, canonical manuscript
README, inspected 10 October 2026.
\href{https://github.com/VladimirReshetnikov/ProveIt/blob/87c7db629552afbc2e803001526debe3430a31b0/Analysis/Polylogarithms/docs/manuscript/README.md}{Repository source anchor}.

\bibitem{repo:dougall}
ProveIt Contributors,
\emph{Mixed Gamma derivatives and exact harmonic tail identities},
\path{chapters/07-mixed-dougall-identities.tex}, in the canonical
Polylogarithms manuscript, source commit
\texttt{87c7db629552afbc2e803001526debe3430a31b0}.
\href{https://github.com/VladimirReshetnikov/ProveIt/blob/87c7db629552afbc2e803001526debe3430a31b0/Analysis/Polylogarithms/docs/manuscript/chapters/07-mixed-dougall-identities.tex}{Pinned chapter source}.

\bibitem{repo:quotients}
ProveIt Contributors,
\emph{Mixed endpoints and all-weight cyclotomic product quotients},
\path{chapters/04-cyclotomic-quotients.tex}, same source commit.
The first 240 source lines were inspected for conventions and the
formal-versus-numerical distinction.
\href{https://github.com/VladimirReshetnikov/ProveIt/blob/87c7db629552afbc2e803001526debe3430a31b0/Analysis/Polylogarithms/docs/manuscript/chapters/04-cyclotomic-quotients.tex}{Pinned chapter source}.

\bibitem{repo:intake}
ProveIt Contributors,
\emph{Incoming reports}, \path{docs/incoming/README.md}, intake procedure
and directory inventory inspected during this research session.
The binary incoming archives were not audited here.
\href{https://github.com/VladimirReshetnikov/ProveIt/tree/main/docs/incoming}{Repository drop zone}.

\bibitem{dlmf:gauss}
NIST Digital Library of Mathematical Functions,
\S15.4, equation 15.4.20: Gauss hypergeometric summation.
\url{https://dlmf.nist.gov/15.4.E20}.

\bibitem{dlmf:gamma}
NIST Digital Library of Mathematical Functions,
\S5.11: asymptotic expansions and ratios of Gamma functions.
\url{https://dlmf.nist.gov/5.11}.

\bibitem{dlmf:beta}
NIST Digital Library of Mathematical Functions,
\S5.12, equation 5.12.1: Euler's beta integral.
\url{https://dlmf.nist.gov/5.12.E1}.

\bibitem{dlmf:polygamma}
NIST Digital Library of Mathematical Functions,
\S5.15: polygamma functions, shifts, and special values.
\url{https://dlmf.nist.gov/5.15}.

\bibitem{dlmf:hurwitz}
NIST Digital Library of Mathematical Functions,
\S25.11: Hurwitz zeta, parameter differentiation, special values,
and the log-Gamma identity.
\url{https://dlmf.nist.gov/25.11}.

\bibitem{dlmf:polylog}
NIST Digital Library of Mathematical Functions,
\S25.12: polylogarithms and the dilogarithm reflection identity.
\url{https://dlmf.nist.gov/25.12}.

\bibitem{krattenthaler}
C.~Krattenthaler,
\emph{Advanced determinant calculus},
S\'eminaire Lotharingien de Combinatoire \textbf{42} (1999), Article B42q,
67 pp.; arXiv:math/9902004.
\url{https://arxiv.org/abs/math/9902004}.
Cited as classical determinant-method background, not as a priority
attribution for every formula in this report.

\end{thebibliography}

% END references.tex

\end{document}
